\chapter{System Design}

\section{System Architecture}
The architecture of \textbf{Event Blinker} follows a three-tier client-server model with real-time communication:

\begin{itemize}[noitemsep]
    \item \textbf{Presentation Layer:} React Native (Expo) mobile app + React/Vite web portal + HTML admin portal.
    \item \textbf{Application Layer:} Node.js/Express.js API server with Socket.io real-time engine.
    \item \textbf{Data Layer:} PostgreSQL with PostGIS extension for geospatial data, with file storage for uploads.
\end{itemize}

% ---- FIGURE: System Architecture ----
% >>> CREATE THIS: A clean 3-tier diagram showing:
%     [Mobile App (Expo)] [Web Portal (React/Vite)] [Admin Portal (HTML)]
%               ↓ HTTP/REST + WebSocket (Socket.io) ↓
%     [Node.js/Express API Server]
%         ├── JWT Auth Middleware
%         ├── 11 Route Modules
%         └── Socket.io Handler
%               ↓
%     [PostgreSQL + PostGIS]     [OpenRouter AI API]    [File Storage]
\begin{figure}[H]
    \centering
    \fbox{\parbox{0.85\textwidth}{\centering\vspace{4cm}
    \textbf{[FIGURE: Three-Tier System Architecture Diagram]}\\[0.3cm]
    \small \textbf{Tier 1 (Clients):} Mobile App (Expo) | Web Portal (React/Vite) | Admin Portal (HTML)\\
    $\downarrow$ REST API + WebSocket $\downarrow$\\
    \textbf{Tier 2 (Server):} Node.js/Express (JWT Auth + 11 Routes + Socket.io)\\
    $\downarrow$\\
    \textbf{Tier 3 (Data):} PostgreSQL/PostGIS | OpenRouter AI | File Storage\\[0.3cm]
    \textit{Create this as a professional block diagram with arrows.}
    \vspace{2cm}}}
    \caption{System Architecture of Event Blinker}
    \label{fig:system_architecture}
\end{figure}

\section{Database Design}

\subsection{Entity-Relationship Overview}
The database consists of 12+ tables organized into three domains:

\begin{table}[H]
\centering
\caption{Database Schema Overview}
\label{tab:db_schema}
\begin{tabular}{@{}llp{6.5cm}@{}}
\toprule
\textbf{Domain} & \textbf{Table} & \textbf{Key Fields} \\
\midrule
\multirow{2}{*}{Users} & \texttt{users} & id (UUID), email, password, user\_type, is\_verified \\
                        & \texttt{organizer\_profiles} & company\_name, verification\_status \\
\midrule
\multirow{4}{*}{Events} & \texttt{events} & location\_geom (PostGIS), category, is\_approved \\
                         & \texttt{event\_likes} & user\_id, event\_id (UNIQUE constraint) \\
                         & \texttt{check\_ins} & user\_id, event\_id, check\_in\_time \\
                         & \texttt{chat\_messages} & sender\_type (user/organizer/bot) \\
\midrule
\multirow{5}{*}{Rides} & \texttt{riders} & current\_location (PostGIS), registration\_status \\
                        & \texttt{vehicles} & vehicle\_type, license\_plate, billbook\_photo \\
                        & \texttt{driver\_licenses} & verification\_status, license\_photo\_url \\
                        & \texttt{ride\_requests} & pickup/dropoff (PostGIS), status, distance\_km \\
                        & \texttt{ride\_offers} & offered\_price, estimated\_arrival \\
\bottomrule
\end{tabular}
\end{table}

% ---- FIGURE: ER Diagram ----
% >>> CREATE THIS: An ER diagram showing relationships between:
%     users (1) → (M) events, users (1) → (M) event_likes, users (1) → (1) riders
%     riders (1) → (1) vehicles, riders (1) → (1) driver_licenses
%     riders (1) → (M) ride_requests, events (1) → (M) chat_messages
\begin{figure}[H]
    \centering
    \fbox{\parbox{0.85\textwidth}{\centering\vspace{4cm}
    \textbf{[FIGURE: Entity-Relationship (ER) Diagram]}\\[0.3cm]
    \small Show tables: users, events, event\_likes, check\_ins, chat\_messages, riders, vehicles, driver\_licenses, ride\_requests, ride\_offers, ride\_reviews\\
    Show foreign key relationships with cardinality (1:M, 1:1)\\[0.3cm]
    \textit{Create this using draw.io or dbdiagram.io}
    \vspace{2cm}}}
    \caption{Entity-Relationship Diagram of Event Blinker Database}
    \label{fig:er_diagram}
\end{figure}

\section{System Diagrams}

\subsection{Use Case Diagram}
The use case diagram identifies four primary actors and their interactions with the system:
\begin{itemize}[noitemsep]
    \item \textbf{User/Attendee:} Discover events on map, view details, like, check-in, chat, book rides, interact with Blinker AI.
    \item \textbf{Event Organizer:} Create events, manage attendees, view analytics, moderate chat.
    \item \textbf{Driver/Rider:} Register with vehicle/license, go online, accept ride requests, offer custom prices.
    \item \textbf{Admin:} Approve/reject organizers, events, riders, and licenses.
\end{itemize}

% ---- FIGURE: Use Case Diagram ----
% >>> CREATE THIS: UML Use Case diagram with 4 actors and their use cases as listed above
\begin{figure}[H]
    \centering
    \fbox{\parbox{0.85\textwidth}{\centering\vspace{4cm}
    \textbf{[FIGURE: Use Case Diagram]}\\[0.3cm]
    \small \textbf{Actors:} User, Organizer, Driver, Admin\\
    \textbf{User Cases:} Discover Events, View Detail, Like, Check-in, Chat, Book Ride, Ask Blinker AI\\
    \textbf{Organizer Cases:} Create Event, View Analytics, Moderate Chat\\
    \textbf{Driver Cases:} Register, Go Online, Accept Rides, Offer Price\\
    \textbf{Admin Cases:} Approve Organizer, Approve Event, Verify Rider, Verify License\\[0.3cm]
    \textit{Create this as a standard UML Use Case diagram.}
    \vspace{2cm}}}
    \caption{Use Case Diagram of Event Blinker}
    \label{fig:usecase}
\end{figure}

\subsection{Sequence Diagram}
The sequence diagram illustrates the event discovery and AI chat flow:

% ---- FIGURE: Sequence Diagram ----
% >>> CREATE THIS: UML Sequence diagram:
%     User → Mobile App → Express API → PostGIS (find nearby events)
%     User → Mobile App → Express API → OpenRouter (AI chat) → Response
%     Express API → Socket.io → Mobile App (real-time updates)
\begin{figure}[H]
    \centering
    \fbox{\parbox{0.85\textwidth}{\centering\vspace{4cm}
    \textbf{[FIGURE: Sequence Diagram -- Event Discovery \& AI Chat]}\\[0.3cm]
    \small \textbf{Lifelines:} User, Mobile App, Express API, PostgreSQL/PostGIS, OpenRouter AI, Socket.io\\
    \textbf{Flow 1:} User opens map $\rightarrow$ App sends GPS coords $\rightarrow$ API runs ST\_DWithin $\rightarrow$ Returns events $\rightarrow$ Markers rendered\\
    \textbf{Flow 2:} User asks AI $\rightarrow$ API fetches event context $\rightarrow$ Sends to LLAMA 3.3 $\rightarrow$ Returns reply\\[0.3cm]
    \textit{Create this as a standard UML Sequence diagram.}
    \vspace{2cm}}}
    \caption{Sequence Diagram: Event Discovery and AI Interaction}
    \label{fig:sequence}
\end{figure}

\subsection{Activity Diagram}
The activity diagram shows the complete ride-sharing booking flow:

% ---- FIGURE: Activity Diagram ----
% >>> CREATE THIS: UML Activity diagram:
%     Start → Passenger sets pickup/dropoff → System calculates distance (Haversine)
%     → System calculates price → Passenger confirms → Socket.io notifies riders
%     → [Decision] Rider accepts OR Rider offers custom price
%     → Passenger confirms → Ride starts → Ride completes → Rating/Review → End
\begin{figure}[H]
    \centering
    \fbox{\parbox{0.85\textwidth}{\centering\vspace{4cm}
    \textbf{[FIGURE: Activity Diagram -- Ride Booking Flow]}\\[0.3cm]
    \small \textbf{Steps:} Start $\rightarrow$ Set Pickup/Dropoff $\rightarrow$ Calculate Distance (Haversine) $\rightarrow$ Calculate Price (NPR) $\rightarrow$ Create Request $\rightarrow$ Notify Riders (Socket.io)\\
    $\rightarrow$ [Decision Diamond] Rider Accepts / Rider Offers Custom Price\\
    $\rightarrow$ Ride In-Progress $\rightarrow$ Ride Complete $\rightarrow$ Review \& Rating $\rightarrow$ End\\[0.3cm]
    \textit{Create this as a standard UML Activity diagram.}
    \vspace{2cm}}}
    \caption{Activity Diagram: Ride Booking Process}
    \label{fig:activity}
\end{figure}
